\chapter{Discussion}


\section{Overview of Main Findings}
The main goal of this thesis was to investigate whether a slice-based CT pipeline could predict patient-level mucus burden from volumetric chest CT scans under weak supervision. The results show that this is possible to some extent, but also that the task remains difficult. The best internal model did not achieve perfect prediction, but the pattern across the experiments suggests that several parts of the pipeline may have contributed useful information. This section summarizes the findings in relation to RQ1--RQ5 before the following sections discuss each component in more detail.

The baseline experiments showed that the representative slice selection performed better than uniform slice sampling. This is an important result because the model does not process the full CT volume directly. It depends on which slices are selected before feature extraction. The comparison suggests that selecting slices based on central restriction, informativeness filtering, and diversity sampling may give the model a better input set than simply sampling slices evenly through the volume.

The main internal experiments also showed better results for the three-channel high-pass representation than for the two-window input, and the strongest internal result was obtained when ORB-BoVW feature fusion was added. The \textbf{ResNet18-3CH+BoVW (Proposed model)} achieved the lowest MAE and RMSE, and the highest \(R^2\) and Spearman correlation among the reported internal configurations. This suggests that the learned CNN representation and the handcrafted local-feature branch may have captured some complementary information.

The backbone comparison did not show that a larger or more domain-specific pretrained backbone was automatically better. DenseNet121 with RadImageNet initialization gave slightly better MAE and RMSE in the three-channel setting, but ResNet18 performed better in the two-window and fusion settings and had stronger rank correlation overall. In this dataset ResNet18 appeared to be the more reliable backbone across different configurations.

The transfer-learning experiments showed that fine-tuning was more useful than zero-shot transfer in both transfer directions. This is expected because the internal mucus dataset and MosMed differ in both dataset properties and target definition. The transfer results therefore support the idea that pretrained or source-trained representations can be useful under domain shift, but they still need target-domain adaptation to perform well.

Overall, the results support the main design direction of the thesis. In the tested configurations, patient-level CT prediction was associated with better performance when slice selection, richer input representation, and feature fusion were included. At the same time, the findings remain exploratory because the internal cohort is small, supervision is only available at patient level, and the transfer setup uses a different external target. These limitations are discussed later, but they are important context for reading the component-wise patterns.

These findings answer the research questions from Chapter~1 in a cautious way. For RQ1, the results show that patient-level mucus-burden prediction from CT is possible to some extent under weak supervision, but still difficult. RQ2 is addressed by the baseline comparison, where the representative slice selection performed better than uniform sampling. RQ3 is addressed by the input-representation and feature-fusion results, where the high-pass and ORB-BoVW configurations were associated with better internal model results. RQ4 is addressed by the backbone comparison, where ResNet18 was more reliable overall than DenseNet121 in this setup. Finally, RQ5 is addressed by the transfer-learning experiments, where fine-tuning was more useful than zero-shot transfer across the two source--target directions. The following sections discuss these points in more detail by linking the results back to the main pipeline components and their limitations.


\section{Ablation-Style Interpretation of Pipeline Components}
The results can be read as an ablation-style progression from simpler settings toward the final pipeline. This section mainly supports RQ2, RQ3, and RQ4 by examining what happens when important components are added or compared, first the slice-selection strategy, then the high-pass input channel, then handcrafted feature fusion, and finally the choice of CNN backbone. These comparisons help show which design choices appeared most useful in the final patient-level prediction setup.

At the same time, the comparisons should not be interpreted as perfect single-factor ablations in every case. Some comparisons in this thesis are ablation-style rather than strict single-factor ablations, because multiple settings may differ between configurations. The results therefore indicate which configurations were associated with better performance, but they do not always isolate the causal effect of each component. This is especially important because the dataset is small and the validation folds contain only a limited number of patients.


\subsection{Representative Slice Selection Strategy}
The first component to consider is the representative slice-selection strategy, which directly relates to RQ2. Since the model does not process the full CT volume directly, the selected slices define the image evidence that the CNN receives. If this step selects many uninformative slices, the later model stages have less useful input, even if the CNN architecture itself is unchanged.

The baseline comparison showed that representative slice selection performed better than uniform sampling. With uniform sampling, the ResNet18 two-window baseline obtained an MAE of \(1.922\), RMSE of \(2.268\), \(R^2\) of \(-0.009\), and Spearman correlation of \(0.202\). When the representative slice selection was used, the corresponding values were an MAE of \(1.760\), RMSE of \(2.223\), \(R^2\) of \(0.003\), and Spearman correlation of \(0.306\). The difference is not dramatic for every metric, but the direction is consistent.

This is consistent with the slice-selection strategy giving the model a better patient-level slice set than simply sampling evenly through the scan. This makes sense for the task, because mucus-related evidence is not expected to be equally visible in every axial slice. Some slices may contain little relevant lung content, while others may contain more useful local information. The central restriction, informativeness filtering, and diversity sampling may therefore reduce the chance that the model spends much of its capacity on irrelevant or weak slices.

Another advantage of the representative strategy is that it reduces dependence on which axial positions happen to be sampled. A uniform strategy may include useful slices, but it can also miss informative regions or include many slices with little relevant lung content, especially when the number of selected slices is limited. This can make the baseline more sensitive to scan length, slice spacing, and the exact sampled positions. In contrast, the representative strategy uses fixed rules for central restriction, informativeness filtering, and diversity sampling, so the selected set is guided by image content rather than by axial position alone.

However, this should be read as patient-level slice selection, not localization. The strategy is heuristic, unsupervised with respect to mucus burden, and does not use slice-level mucus labels. Its role is more modest: it creates a more useful and stable subset of slices for patient-level prediction. Within that role, the baseline comparison supports keeping the representative slice-selection stage in the final pipeline.


\subsection{Input Representation and High-Pass Channel}
The next step in the pipeline was the move from the two-window input to the three-channel input, which is part of RQ3. The two-window representation already gives the model two useful views of the CT slice, one focused on lung-window information and one focused on mediastinal-window information. The three-channel version keeps these two channels and adds a high-pass channel. The purpose of the high-pass channel is to emphasize local intensity changes and fine image structure.

The internal results indicate that the three-channel input outperformed the two-window input. The ResNet18 two-window model achieved a mean absolute error (MAE) of \(1.451\), root mean square error (RMSE) of \(1.858\), coefficient of determination (\(R^2\)) of \(0.126\), and Spearman correlation of \(0.363\). With the high-pass third channel, the ResNet18 three-channel model used the same main pipeline configuration and achieved an MAE of \(1.386\), RMSE of \(1.826\), \(R^2\) of \(0.186\), and Spearman correlation of \(0.387\). Although the improvement is smaller than the difference observed with \textbf{ResNet18-3CH+BoVW (Proposed model)}, it remains consistent across all reported metrics.

This suggests that the high-pass channel may have added useful information beyond the original two CT windows. A possible explanation is that the high-pass channel makes local changes in the image more visible to the CNN. This can be useful in a task where the relevant evidence may be small, local, and not evenly distributed across the lungs. The implementation of lung-mask suppression also helps focus this third-channel response inside the approximate lung region instead of letting high-frequency structures outside the lungs dominate the input.

At the same time, the high-pass channel should not be interpreted as a direct mucus detector. It highlights local image structure rather than clinically verified mucus findings. The safer interpretation is that it gave the model an additional view of the selected slices, and that this extra view was useful for patient-level prediction in the reported experiments.


\subsection{Handcrafted ORB-BoVW Feature Fusion}
The strongest internal result for RQ3 was obtained when the CNN features were combined with the ORB-BoVW handcrafted feature branch. Compared with the three-channel CNN-only model, the \textbf{ResNet18-3CH+BoVW (Proposed model)} reduced the MAE from \(1.386\) to \(\mathbf{1.244}\), reduced the RMSE from \(1.826\) to \(\mathbf{1.712}\), and increased the Spearman correlation from \(0.387\) to \(\mathbf{0.498}\). This was the clearest performance gain associated with a configuration change in the internal experiments, and it suggests that the handcrafted branch may have added information that the CNN did not fully capture on its own.

One possible reason is that ORB-BoVW describes the image in a different way than the CNN. The CNN learns features directly from the training data, while the ORB-BoVW branch summarizes local keypoint patterns using a fixed visual vocabulary. In a larger dataset the CNN might learn many of these local structures by itself. In this thesis the internal cohort is small, so the handcrafted branch may have helped by adding a more stable description of local image texture and structure. This fits the general idea behind the fusion setup, the handcrafted features are not meant to replace the CNN, but to give the model another type of signal.

The result is also interesting because mucus-burden prediction is likely influenced by local and unevenly distributed findings. A patient-level score can depend on small regions in a subset of slices rather than on the whole scan equally. The high-pass channel already gives the CNN an input that emphasizes local intensity changes, but the ORB-BoVW branch adds another summary of local image patterns. The fusion result suggests that these two views were complementary in this experimental setting.

At the same time, ORB-BoVW keypoints are not mucus-specific. The branch describes local image structure through keypoints and visual words, so its benefit should be interpreted as additional patient-level signal rather than direct mucus detection.

There is also a trade-off. Adding ORB-BoVW fusion increases the pipeline's complexity. It requires fitting a visual vocabulary within each training fold, caching or computing handcrafted features, and making sure validation data are not used when building the vocabulary. This increases the number of moving parts compared with a CNN-only model. Because the fusion configuration had better values for all reported internal metrics, the extra complexity appears justified in this thesis setup.

Overall, the ORB-BoVW result supports the idea that handcrafted features can still be useful in small medical-imaging datasets. The best interpretation is not that handcrafted features are better than learned features, but that they can complement learned CNN features when the amount of training data is limited and the target is only available at patient level.


\subsection{Backbone Comparison}
The backbone comparison addresses RQ4 by comparing whether a different CNN feature extractor was associated with better patient-level prediction. In this thesis, ResNet18 and DenseNet121 were compared across the same main settings, the two-window input, the three-channel input, and the ORB-BoVW fusion. DenseNet121 is a deeper architecture and was initialized with RadImageNet weights, so it was reasonable to expect that it might provide stronger medical-image features. However, the results did not show a clear overall advantage from switching to DenseNet121.

For the two-window input, ResNet18 performed better than DenseNet121 across all reported metrics. ResNet18 achieved an MAE of \(1.451\), RMSE of \(1.858\), and Spearman correlation of \(0.363\), while DenseNet121 obtained an MAE of \(1.602\), RMSE of \(2.038\), and Spearman correlation of \(0.208\). This suggests that the larger backbone was not beneficial in the simpler two-window setting.

The three-channel setting was more mixed. DenseNet121 gave slightly lower MAE and RMSE than ResNet18, with MAE decreasing from \(1.386\) to \(1.356\) and RMSE from \(1.826\) to \(1.789\). However, ResNet18 still had better \(R^2\) and Spearman correlation. This means that DenseNet121 reduced the average error slightly in this setting, but it did not rank patients as well as ResNet18. Since Spearman correlation is important for understanding whether the model orders patients correctly by predicted burden, this difference matters.

The fusion setting again favoured ResNet18. \textbf{ResNet18-3CH+BoVW (Proposed model)} achieved the best overall internal result, while DenseNet121 performed worse on MAE, RMSE, \(R^2\), and Spearman. If DenseNet121 had been clearly better as a backbone in this setup, one would expect it to perform well in this setting as well, but that was not observed.

One possible explanation is that DenseNet121 may have been harder to adapt in this small dataset. The internal cohort contains only 32 patient cases, and each validation fold contains only a few patients. In such a setting, a larger or more complex backbone does not automatically lead to better generalization. ResNet18 is simpler, and this may have made it more stable when combined with the selected slice representation, aggregation strategy, and handcrafted fusion branch.

This does not mean that ResNet18 is generally better than DenseNet121 for CT analysis. It only means that, in this experimental setup, DenseNet121 did not provide a consistent advantage. The results suggest that backbone choice should be evaluated together with the rest of the pipeline, rather than assuming that a deeper or radiology-pretrained model will always perform better. For this thesis, ResNet18 was therefore kept as the main backbone because it gave the strongest and most consistent results across the reported configurations.


\subsection{Aggregation and Patient-Level Prediction}
Aggregation was an important part of the pipeline because the model produces slice-level scores, while the target is defined at patient level. During development, this was approached step by step. The simplest option was mean aggregation, where all selected slice scores contribute equally to the final patient prediction. This is stable and easy to interpret, but it assumes that all selected slices are equally informative.

The next option was top-\(k\) aggregation. This was motivated by the idea that mucus-related evidence may only be visible in a subset of slices. Instead of averaging all slice scores, top-\(k\) aggregation focuses on the highest predicted slice scores. This can be useful when the patient-level target is driven by the most abnormal or informative slices, but it also risks ignoring broader information from the rest of the scan.

The final hybrid aggregation strategy was built from these two ideas. It combines the mean prediction with the top-\(k\) prediction, so the final patient-level estimate keeps information from all selected slices while still giving extra weight to the slices with stronger predicted evidence. This made sense for the task because mucus burden is patient-level, but the visual evidence can be local and unevenly distributed.

This aggregation step also reflects the weakly supervised setup. The slice-level scores are useful only after being combined into a patient-level prediction; they are not independently verified slice findings. The purpose of aggregation is therefore to connect these intermediate scores to the available patient-level target.

Although the final pipeline uses hybrid aggregation, this thesis does not present a separate result table comparing mean, top-\(k\), and hybrid aggregation. Therefore, the discussion should not claim that hybrid aggregation was proven to be optimal. A more careful interpretation is that hybrid aggregation was chosen because it combines the stability of mean aggregation with the more focused behaviour of top-\(k\) aggregation. A full aggregation ablation would be a useful direction for future work.


\subsection{Target Transformation and Reporting Scale}
The target scale also affects how the model is trained and how the results should be interpreted. In the baseline experiments, the raw mucus-burden score was used directly. In the main internal experiments and transfer-learning experiments, the models were trained using the log-transformed target, and predictions were transformed back to the raw mucus-score scale for reporting.

This choice was made because the raw mucus-burden scores are not evenly distributed. Most patients have relatively low or moderate values, while a few patients have higher scores. In this situation, training directly on the raw scale can make the model more sensitive to larger target values. The log transformation can reduce the influence of these larger values and may make the regression problem easier to optimize.

At the same time, reporting MAE and RMSE on the log scale would be harder to interpret clinically and practically. For that reason, predictions were transformed back to the raw score scale before reporting the main error metrics. This makes the reported MAE and RMSE easier to understand, because they are expressed in the same units as the original mucus-burden score.

The log transformation should also be interpreted carefully in the ablation argument. The raw-target baseline and the log-target main experiments do not differ only in target scale, they also differ in other parts of the configuration. Therefore, the results do not isolate the effect of the log transformation alone. It is more accurate to say that the log target was part of the final training setup that gave the best internal results, rather than claiming that the log transformation alone caused the improvement.

The target transformation was part of a training setup intended to handle a small and unevenly distributed regression target, while the back-transformation kept the reported results understandable on the original score scale. This balance is important because the model should be optimized in a stable way, but the results still need to be interpretable in terms of patient-level mucus burden.


\section{Transfer Learning and Domain Shift}
The transfer-learning experiments address RQ5 by asking whether external CT data with a different target can support representation transfer in a small-cohort setting. This matters because the internal mucus dataset is small, so it is natural to ask whether a model trained on another CT task can provide a useful starting point.

The transfer experiments used two model variants from the main pipeline. The first was ResNet18-3CH, which refers to the ResNet18 model with the main three-channel high-pass input. The second was \textbf{ResNet18-3CH+BoVW (Proposed model)}, which adds the ORB-BoVW handcrafted feature-fusion. They include the main configuration for the slice-based pipeline, input representation, and patient-level aggregation described in Chapter~4.

The main pattern is clear, fine-tuning was more useful than zero-shot transfer in both directions. This suggests that the models could reuse some information from the source domain, but they still needed target-domain adaptation. This is expected because the internal mucus dataset and MosMed are both based on chest CT, but they do not have the same target. The internal task predicts mucus burden, while MosMed uses a case-level CT severity label related to COVID-19 lung involvement.


\subsection{Zero-Shot Transfer}
Zero-shot transfer means that a model trained on the source domain is evaluated directly on the target domain without any target-domain training. In this thesis, zero-shot transfer was limited in both directions.

In the mucus-to-MosMed direction, the ResNet18-3CH zero-shot model obtained an MAE of \(0.903\), RMSE of \(1.070\), \(R^2\) of \(-1.62\), and Spearman correlation of \(-0.027\). The \textbf{ResNet18-3CH+BoVW (Proposed model)} zero-shot model was slightly better on MAE and RMSE with values of \(\mathbf{0.856}\) and \(\mathbf{1.013}\), but it still had a negative \(R^2\) of \(\mathbf{-0.95}\) and a negative Spearman correlation of \(\mathbf{-0.160}\). This shows that even the proposed model trained on the mucus task did not transfer well to the MosMed severity task without adaptation.


The MosMed-to-mucus direction showed a similar pattern. The ResNet18-3CH zero-shot model obtained an MAE of \(1.889\), RMSE of \(2.199\), \(R^2\) of \(0.08\), and Spearman correlation of \(0.050\). The \textbf{ResNet18-3CH+BoVW (Proposed model)} zero-shot model performed better, with an MAE of \(\mathbf{1.802}\), RMSE of \(\mathbf{2.097}\), \(R^2\) of \(\mathbf{0.10}\), and Spearman correlation of \(\mathbf{0.261}\). This suggests that the BoVW fusion branch transferred somewhat better than the non-fusion variant in this direction, but the zero-shot setting was still weaker than fine-tuning.

These results are not surprising. Zero-shot transfer assumes that the source-domain model can be applied directly to a different dataset and a different target definition. Even though both datasets contain chest CT scans, the model is not solving the same prediction task in both cases. The weak zero-shot results therefore support the idea that representation reuse alone is not enough for this problem.


\subsection{Fine-Tuning}
The fine-tuned setting performed better than zero-shot transfer in both directions. This means that the source-domain model was more useful as an initialization than as a final model applied directly to the target domain.

In the mucus-to-MosMed direction, the fine-tuned ResNet18-3CH model had a lower MAE than the zero-shot setting, from \(0.903\) to \(0.460\), and a lower RMSE, from \(1.070\) to \(0.643\). The \textbf{ResNet18-3CH+BoVW (Proposed model)} showed the same pattern, with MAE changing from \(0.856\) to \(\mathbf{0.369}\), and RMSE from \(1.013\) to \(\mathbf{0.545}\). The fine-tuned \textbf{ResNet18-3CH+BoVW (Proposed model)} also achieved the strongest MosMed-target result, with \(R^2=\mathbf{0.29}\) and Spearman correlation \(\mathbf{0.489}\).


In the MosMed-to-mucus direction, the fine-tuned settings also had better error metrics than the zero-shot settings. The ResNet18-3CH model changed from an MAE of \(1.889\) to \(1.287\), and from an RMSE of \(2.199\) to \(1.576\). The \textbf{ResNet18-3CH+BoVW (Proposed model)} changed from an MAE of \(1.802\) to \(\mathbf{1.329}\), and from an RMSE of \(2.097\) to \(\mathbf{1.605}\). In this direction, the fine-tuned ResNet18-3CH model gave the lowest MAE, lowest RMSE, and highest \(R^2\), while the fine-tuned \textbf{ResNet18-3CH+BoVW (Proposed model)} gave the highest Spearman correlation.


The main interpretation is that source-domain training can provide a useful starting point, but the model still needs to adapt to the target dataset. This is especially relevant in small medical-imaging datasets, where training from scratch is difficult, but direct zero-shot transfer is usually too weak. Fine-tuning gives a middle ground, the model starts from features learned on another CT task, and then the later training adjusts those features to the target-domain label.


\subsection{Transfer Direction and Dataset Mismatch}
The transfer results should be interpreted with the dataset mismatch in mind. MosMed is not an external mucus-burden validation dataset. Its target is a different case-level CT severity label, while the internal target is a mucus-burden score. This means that the transfer experiments test whether CT representations can be adapted between related tasks, and not whether the mucus model generalizes directly to an external mucus-scored cohort.

This target mismatch helps explain why fine-tuning was needed. A model trained on mucus burden may learn useful CT features, but those features still need to be connected to the MosMed severity label. The same is true in the opposite direction, a MosMed-trained model may learn features related to lung involvement, but it must be adapted before it can predict mucus burden.

The direction of transfer also matters. In the mucus-to-MosMed direction, the fine-tuned settings performed better than the zero-shot settings for both ResNet18-3CH and \textbf{ResNet18-3CH+BoVW (Proposed model)}. In this direction, the proposed model was strongest across all reported metrics. In the MosMed-to-mucus direction, the fine-tuned settings also performed better than the zero-shot settings, but the best model depended on the metric. ResNet18-3CH had the lowest MAE, lowest RMSE, and highest \(R^2\), while \textbf{ResNet18-3CH+BoVW (Proposed model)} had the highest Spearman correlation. This suggests that transfer behaviour is not fully symmetric between the two datasets.

Overall, the transfer experiments answer the transfer-related research question in a cautious way. They show that external CT data and source-domain training can be useful for representation transfer, but mainly when followed by target-domain fine-tuning. They also show that zero-shot transfer is not enough when the dataset and target definition change. For this thesis, transfer learning should therefore be viewed as a useful adaptation strategy, not as external clinical validation of mucus-burden prediction.



\section{Weak Supervision and Interpretation of Slice-Level Evidence}
A central challenge in this thesis is that the available labels are defined at patient level, while the model is trained on selected CT slices. This point is central to RQ1 and RQ2, because the thesis evaluates patient-level prediction while also asking how the selected slice set affects that prediction. During training, each selected slice from patient \(i\) is paired with the same patient-level target \(y_i\). For example, if a patient has a mucus-burden score of \(5\), then all selected slices from that patient are trained with target value \(5\), even though the actual mucus evidence may only be visible in some of those slices.

This is the weakly supervised part of the setup. The label is correct for the patient, but it is not a verified label for each individual slice. A slice-level score \(s_{i,j}\) may be useful for the final patient-level prediction, but it should not be treated as a true slice-level mucus measurement.

At validation and testing time, the model produces slice-level scores for the selected slices. These scores are aggregated into one patient-level prediction \(\hat{y}_i\), which is compared with the true patient-level target \(y_i\). The evaluation is therefore patient-level, not slice-level. The pipeline can be judged by how well it predicts the patient score, but not by whether individual slice scores correctly localize mucus.

This also affects how the slice-selection strategy and the high-pass input representation should be understood. These components may help the model build a better patient-level prediction, but they should not be read as clinical explanation maps. Different slice-score patterns could also lead to similar patient-level predictions; a patient score could be driven by a few high-scoring slices or by many moderately scoring slices. Without slice-level annotations, it is not possible to know which pattern is more clinically correct.

For this reason, the results should be interpreted as patient-level prediction results, not localization results. The pipeline suggests that selected slices and local image features contain useful information for predicting mucus burden, but it does not prove that the model has learned to identify mucus plugs directly. A future version of the project would need slice-level or region-level annotations to evaluate that question properly.


\section{Reliability, Generalization, and Limitations}
The main limitation of this thesis is the small size of the internal dataset. The internal experiments were performed on 32 patient cases, with only six or seven patients in each validation fold. This makes the results sensitive to individual cases. A single difficult or unusual patient can have a large effect on fold-level metrics, especially for RMSE and \(R^2\). The standard deviations reported in Chapter~5 should therefore be taken seriously, not treated as small details.

The low and sometimes negative \(R^2\) values are also important for interpreting the results honestly. They show that the models do not consistently explain a large share of target variance. Some settings perform close to, or worse than, a simple mean-prediction reference. This is not unexpected given the small internal cohort, weak patient-level supervision, possible noise in the mucus-score target, and the difficulty of mapping selected CT slices to one patient-level burden score. MAE and RMSE are therefore important complementary metrics because they show the size of the prediction error on the target scale, while Spearman correlation helps indicate whether patients are ranked in a meaningful order.

The small cohort also limits how strongly the ablation-style findings can be generalized. The results suggest that representative slice selection strategy, high-pass input representation, and ORB-BoVW fusion were useful in this dataset. However, the same performance patterns may not appear in the same way on a larger or more diverse cohort. This is especially relevant for the handcrafted fusion branch. It performed well here, but it also adds more preprocessing and feature-engineering steps, which may behave differently on data from other scanners, hospitals, or patient groups.

The patient-level supervision also limits interpretability. Even when a model predicts the patient score reasonably, it cannot be checked directly against slice-level mucus findings in this dataset. This is the main reason the thesis avoids localization claims.

The transfer-learning experiments add a second generalization limit. MosMed is useful for studying transfer between related CT tasks, but its target is not mucus burden. These results therefore describe domain-shift behaviour, not external mucus-specific validation. A true external validation would require another dataset with the same or a closely comparable mucus-burden target.

The negative \(R^2\) values in the zero-shot transfer experiments are especially relevant here. They indicate that direct source-to-target transfer can be worse than predicting the target-domain mean when the dataset and label definition change. This reinforces the main transfer-learning interpretation that source-trained CT representations may be useful after fine-tuning, but direct zero-shot reuse is not reliable in this setting.

There are also limitations in the experimental coverage. The thesis compares several important components, but it does not provide a complete ablation of every hyperparameter. For example, the final aggregation strategy was motivated by development experiments with mean and top-\(k\) aggregation, but Chapter~5 does not report a separate aggregation comparison table. Similarly, several third-channel variants were explored during development, including high-pass, DoG, LoG, and DCT-based representations. The high-pass representation was kept as the main three-channel setup because it gave the strongest and most stable results in these development runs. However, these exploratory comparisons were not reported as full formal result tables in Chapter~5, since the results chapter focuses on the main pipeline configurations. This means that the thesis supports the high-pass choice as a development-driven design decision, but does not present a complete formal ablation of every third-channel variant. 

The pipeline is based on two-dimensional slice processing rather than full three-dimensional volume modelling. This made the task more manageable with the available data and computational resources, but it also means that the model may miss some spatial context across neighbouring slices. The representative slice-selection strategy and aggregation stages partly address this by selecting multiple slices and combining their predictions, but they are not the same as modelling the full 3D structure of the CT volume.

Overall, the results are promising, but they should be viewed as preliminary. The thesis shows that the proposed pipeline can learn useful patient-level signals from CT under weak supervision, and that several design choices were associated with better performance within the available data. At the same time, stronger conclusions would require a larger cohort, external mucus-specific validation, and more detailed annotation for evaluating slice-level or regional evidence.


\section{Future Work}
Several directions could make the pipeline stronger in future work. The most important one is to evaluate the method on a larger internal cohort. The current dataset contains only 32 patient cases, which limits how confidently the results can be generalized. A larger dataset would make the cross-validation results more stable and would also make it easier to test whether the performance patterns associated with slice selection, high-pass input, and ORB-BoVW fusion still hold when more variation is present. 

A second important direction is external validation on another dataset with mucus-burden labels. Since MosMed was used for transfer and domain-shift analysis, it cannot serve this role. A true external validation would require an independent cohort where mucus burden is scored in a similar way to the internal dataset.

More detailed annotations would also be valuable. Slice-level or region-level mucus annotations would make it possible to test whether high-scoring slices correspond to mucus-related areas, rather than only supporting the final patient-level score.

The aggregation strategy could also be extended in future work. The final pipeline used a hybrid aggregation rule based on mean and top-\(k\) aggregation, which was chosen because it gave a good balance between using information from all selected slices and emphasizing the most informative slice predictions. Future work could present a more detailed comparison of aggregation choices, including mean, max, top-\(k\), hybrid aggregation, attention-based aggregation, and multiple-instance learning approaches. This would make the aggregation part of the pipeline more transparent and help show when each strategy is most useful.

The input representation could also be studied further. The high-pass channel was selected as the main third-channel configuration because it performed best during development and gave a simple way to emphasize local image structure. Other variants, such as DoG, LoG, and DCT-based filtering, were explored but not reported as full result tables in Chapter~5. Future work could present these comparisons more formally, together with the role of lung-mask suppression, the choice of high-pass source channel, and the high-pass scale. This would make the input-representation design more complete and easier to compare.

The ORB-BoVW fusion was associated with the strongest main internal results, but it adds extra complexity to the pipeline. Future work could test different vocabulary sizes, keypoint sources, descriptor types, or simpler handcrafted alternatives. It would also be useful to investigate whether the fusion branch remains helpful when the dataset becomes larger, since a CNN may be able to learn more local structure by itself when more training data are available.

Finally, future work could explore models that use more spatial context. The current pipeline is based on selected two-dimensional slices, which makes the problem manageable with limited data and computation. However, mucus burden is a volumetric finding, and neighbouring slices may contain useful context. A future model could therefore use 2.5D inputs, short slice sequences, recurrent aggregation, attention-based slice pooling, or full 3D CNNs if enough data are available. These approaches could make greater use of volumetric information while still keeping the patient-level prediction task central.


\section{Concluding Discussion}
This thesis investigated a weakly supervised, slice-based pipeline for predicting patient-level mucus burden from chest CT. The results show that the task is possible to approach with the available data, but also that it remains challenging. The best internal model did not produce perfect predictions, but the experimental progression showed that several richer configurations performed better than simpler settings.

The most important internal finding was that the \textbf{ResNet18-3CH+BoVW (Proposed model)} performed best among the reported configurations. The representative slice selection performed better than uniform sampling, the high-pass three-channel configuration performed better than the two-window input, and adding ORB-BoVW feature fusion gave the strongest internal result. These findings are consistent with the idea that the problem depends not only on the CNN backbone itself, but also on how the CT volume is sampled, represented, and combined with complementary local features.

The transfer-learning experiments also gave a useful result. Zero-shot transfer was limited, while the fine-tuned settings performed better in both transfer directions. This suggests that source-domain CT representations can be useful, but they need to be adapted to the target dataset and target label. This is adaptation evidence, not external mucus validation.

A key point throughout the thesis is that the work is weakly supervised. The slice-level scores are intermediate outputs trained from patient-level labels, so the pipeline should be understood as patient-level prediction rather than mucus localization.

Overall, the thesis shows that a carefully designed slice-based CT pipeline can learn useful patient-level signals for mucus-burden prediction in a small-cohort setting. The results support the use of representative slice selection, high-pass input representation, and handcrafted feature fusion as meaningful parts of the pipeline. However, the findings are still preliminary because of the small dataset, weak supervision, and lack of external mucus-specific validation. The method is therefore best viewed as a promising experimental framework that can be strengthened through larger datasets, more detailed annotations, and broader validation.
